\documentclass[11pt]{article}
\usepackage[a4paper,margin=25mm]{geometry}
\usepackage{amsmath,amssymb,amsthm,booktabs,longtable,array,xcolor}
\usepackage[colorlinks=true,linkcolor=blue!40!black,citecolor=blue!40!black,urlcolor=blue!40!black]{hyperref}
\setlength{\parindent}{0pt}
\setlength{\parskip}{5pt plus 1pt}
\setlength{\emergencystretch}{3em}
\allowdisplaybreaks[2]
\newcommand{\Q}{\mathbb Q}
\newcommand{\Z}{\mathbb Z}
\newcommand{\R}{\mathbb R}
\newcommand{\C}{\mathbb C}
\newcommand{\N}{\mathbb N}
\newcommand{\dd}{\,d}
\newcommand{\Hfive}[1]{H_{#1}^{(5)}}
\newcommand{\source}[1]{{\small\color{blue!40!black}#1}\par}
\newcommand{\heading}[2]{\section{#1}\source{#2}}
\newcommand{\vG}{v_p^{\mathrm G}}
\DeclareMathOperator{\ord}{ord}
\newtheorem{theorem}{Theorem}[section]
\newtheorem{proposition}[theorem]{Proposition}
\newtheorem{lemma}[theorem]{Lemma}
\newtheorem{corollary}[theorem]{Corollary}
\theoremstyle{definition}
\newtheorem{definition}[theorem]{Definition}
\newtheorem{remark}[theorem]{Remark}
\numberwithin{equation}{section}
\setcounter{tocdepth}{1}
\hypersetup{pdftitle={The determinant argument for zeta(5): a concise exposition},
 pdfauthor={Expository draft prepared with OpenAI Codex},
 pdfsubject={Fauzan's construction, guided by Persiflage, with the revised arithmetic of Firsching's formalization}}
\begin{document}
\hypersetup{pageanchor=false}
\begin{titlepage}
\vspace*{15mm}
{\LARGE The determinant argument\\[4pt]
for the irrationality of \(\zeta(5)\)\par}
\vspace{8mm}
{\large A concise exposition guided by Persiflage\par}
\vspace{4mm}
Draft prepared with OpenAI Codex\\
14 November 2026
\vspace{9mm}

The proof combines two descriptions of one determinant. Over the real
numbers it is a positive multiple integral, and logarithmic energy makes
it small. At each prime, carefully chosen polynomial bases expose
divisibility shared by its coefficients. The real decay exceeds the
arithmetic cost. A hypothetical rational value of \(\zeta(5)\) would then
give a positive integer smaller than one.

\vspace{4mm}
\textbf{Intellectual sources.}
Frank Calegari's Persiflage post \cite{persiflage} guides the presentation
and its connection to earlier approximation theory. The positive measure
and its Stieltjes transform are due to Marc Pr\'evost
\cite{prevost}; the particular determinant construction and comparison
measure are those of Aabir Fauzan \cite{fauzan}. The local arithmetic and
normalization follow the revised argument in Moritz Firsching's Lean
development \cite{firsching}. Chun Wing Lai's later exposition
\cite{lai} and the earlier drafts in this chat were used to check and
expand the opening identities. These notes claim no new irrationality
theorem or new determinant construction.

\vspace{3mm}
\textbf{Version followed.}
We use \(K=40n,\ N=3n,\ h=37n\), with cutoff \(M=400\) and the full
outer-range rank loss. This is the version documented by Firsching's
development, pinned to commit \texttt{7fe7367}. Persiflage's concluding
\(K=12n,\ N=n\) variant is not the quantitative construction used here.
The arithmetic estimate is derived for the normalization actually defined
below; it is not Fauzan's original \(M=200\) assertion assumed again.

\vfill
The main text presents the construction and the chain of estimates.
A separate detailed companion supplies the local lemmas, error bounds,
and finite tables. The source package includes both versions, the
verification scripts, and the attribution ledger.
\end{titlepage}
\hypersetup{pageanchor=true}
\pagenumbering{roman}
\tableofcontents
\clearpage
\pagenumbering{arabic}
\section{The argument and its sources}\label{sec:orientation}

Write \(\xi=\zeta(5)=\sum_{j\geq1}j^{-5}\).
The useful starting point in Persiflage's account is a positive measure
whose polynomial moments are rational, while certain values of its
Stieltjes transform contain \(\xi\). Multiplying the measure by a
particular rational function then produces a matrix with two simultaneous
descriptions: its entries are affine expressions in \(\xi\), and it is
a Gram matrix. We follow that organizing perspective
\cite[The formulas of Euler and Hermite; Hankel Determinants]{persiflage}.

The historical distinction matters. Pr\'evost's Theorems 1--2
\cite{prevost} supply the underlying transform and moments; his earlier
work on Pad\'e approximation, and subsequent work with Rivoal
\cite{prevost1996,prevost-rivoal}, explain their role in rational
approximation. Rationally modifying an orthogonality measure belongs to
the theory developed by Uvarov, with a determinant formula proved by
Krattenthaler \cite{uvarov,krattenthaler}. Zudilin's determinantal irrationality
criterion \cite{zudilin} is a direct methodological antecedent.
Brown's integral pairings and coefficient-normalization framework
\cite[Sections 3 and 7.4--7.5]{brown} provide further context.
We do not need these general determinant formulas as unproved steps:
the particular identities used here are established below.

The extra work for this construction is quantitative. The numerator
of the rational modifier improves divisibility, but also changes the
real integral. Both effects must be estimated for the same determinant.
Fauzan makes this choice explicit \cite[Sections 2--6]{fauzan}.
For the local estimates we use Firsching's revised finite
partial-fraction argument \cite{firsching}. In particular, the correction
of bounded rank is charged its full rank, and the prime cutoff is 400.
No inequality with the smaller, disputed rank loss is needed.

\subsection{What the estimates must accomplish}

The desired polynomials have integer coefficients and satisfy
\begin{equation}\label{eq:goal}
 \deg Q_n\leq37n,\qquad
 0<Q_n(\xi)\leq e^{-cn^2}
\end{equation}
for a fixed \(c>0\) and all sufficiently large \(n\).
This form of the final contradiction is used by Fauzan
\cite[Section 1.1]{fauzan}; compare Zudilin's determinant argument
\cite[Section 2]{zudilin}.

Suppose \(\xi=a/b\), where \(a\in\Z\) and \(b\in\Z_{>0}\).
Writing \(Q_n(X)=\sum_{j=0}^{37n}q_jX^j\), with zero coefficients added
if necessary, gives
\[
 b^{37n}Q_n(a/b)=\sum_{j=0}^{37n}q_j a^j b^{37n-j}\in\Z.
\]
It is positive by \eqref{eq:goal}, whereas
\[
 0<b^{37n}Q_n(a/b)
 \leq\exp(37n\log b-cn^2)\longrightarrow0.
\]
The denominator \(b\) is fixed. A positive integer cannot tend to zero.
This proves irrationality once \eqref{eq:goal} has been established.

We construct a rational polynomial \(F_K\), then a positive rational
multiplier \(m_K\) such that \(Q_n=m_KF_K\) for \(K=40n\).
Positivity will follow before any size estimate is used.
The analytic calculation bounds \(\log F_K(\xi)\) by a negative
constant times \(K^2\), up to lower-order terms.
The arithmetic calculation bounds \(\log m_K\) by a smaller positive
constant times \(K^2\). Their sum is negative.

\subsection{Where the details live}

The integral representation is derived in full. Sections 5--7 then keep
the main inequalities and the reasons behind them together, so that
the real decay can be compared directly with the arithmetic cost.
The supporting calculations are in the separate
\emph{detailed companion}, supplied as
\texttt{zeta5\_detailed\_companion.pdf}: its Section 5 proves the analytic
bounds, Sections 6--7 prove the local and global arithmetic bounds,
and Appendix C gives the exact arithmetic table.
Those references supply proofs of the stated estimates, not additional
assumptions. The unchanged finite certificates and source ledger accompany
both versions. No formal verification of this prose or local rebuild of
the complete Lean development is claimed.

\section{A rational functional and a positive measure}\label{sec:functional}
\source{Construction: Fauzan \cite[(2.2)--(2.3)]{fauzan} and Lai
\cite[Definition 4.2]{lai}. The presentation follows the functional
viewpoint in Persiflage \cite{persiflage}.}
\subsection{Defining the functional carefully}
The notation involves three different variables. The letter $t$ is the
argument of a rational function. The letter $X$ is a parameter that will
eventually be set equal to $\xi$. The letter $y$ will be an integration
variable. Keeping these roles separate makes the later matrix calculation
much easier to follow.

Let $\mathcal R$ be the vector space over $\Q$ of rational functions in
$\Q(t)$ whose finite poles, if any, are simple and belong to
$\{-1^2,-2^2,-3^2,\ldots\}$. Every $R\in\mathcal R$ has a unique expression
\begin{equation}\label{eq:partial}
 R(t)=\sum_{e=0}^{d}p_et^e+\sum_{j\in J}\frac{c_j}{t+j^2},
 \qquad p_e,c_j\in\Q,
\end{equation}
where $J$ is finite. Polynomial division gives the polynomial part.
At a pole, its coefficient is
$c_j=\left.(t+j^2)R(t)\right|_{t=-j^2}$.
Subtracting these simple fractions leaves a rational function with no finite
poles, hence a polynomial. These observations also prove uniqueness
(with zero terms omitted).

Define the Bernoulli numbers by
\[
 \frac{u}{e^u-1}=\sum_{k\geq0}B_k\frac{u^k}{k!}.
\]
Multiplication by the exponential series determines the $B_k$ recursively
as rational numbers. In particular $B_0=1$, $B_1=-1/2$, $B_2=1/6$,
and $B_4=-1/30$. Also put
\[
 \Hfive{j}=\sum_{v=1}^j v^{-5},\qquad \Hfive{0}=0.
\]
Prescribe a $\Q$-linear functional $\mu_X:\mathcal R\to\Q[X]$ by
\begin{align}
 \mu_X(t^e)&=
 \frac{(-1)^eB_{2e+2}(2e+3)(2e+4)(2e+5)}{24},
 &&e\in\Z_{\geq0},\label{eq:mu-moment}\\
 \mu_X\!\left(\frac1{t+j^2}\right)&=
 j^4\bigl(X-\Hfive{j}\bigr)-\frac14+\frac1{2j},
 &&j\in\Z_{>0}.\label{eq:mu-pole}
\end{align}
For a general $R$, apply these formulas to its decomposition and add with
the coefficients $p_e,c_j$. This is the meaning of extending by linearity.
It gives a well-defined affine polynomial in $X$: only the simple
fractions contribute a term involving $X$.

For example,
\[
 \mu_X(1)=\frac5{12},\qquad
 \mu_X\!\left(\frac1{t+1}\right)=X-\frac34.
\]
Since $t/(t+1)=1-1/(t+1)$, linearity gives
$\mu_X(t/(t+1))=7/6-X$. There is no rule allowing us to multiply
functional values; $\mu_X$ is linear, not multiplicative.

Finally, $\mu_\xi(R)$ means evaluate the polynomial $\mu_X(R)$ at $X=\xi$.
It does not mean substitute $\xi$ for $t$ in the rational function.

\subsection{Why the polynomial values are moments}
\source{Pr\'evost \cite[Theorems 1--2]{prevost};
Fauzan \cite[Proposition 2.2]{fauzan}; Lai \cite[Lemma 4.5]{lai}.}
Define, for $y>0$,
\begin{equation}\label{eq:weight}
 w(y)=\frac{(2\pi)^4y^5}{12}
       \sum_{\ell\geq1}\ell^4e^{-2\pi\ell y}.
\end{equation}
Every summand is positive. The series converges for each $y>0$, so
$w(y)>0$. We will prove
\begin{equation}\label{eq:representation}
 \boxed{\mu_\xi(R)=\int_0^\infty R(y^2)w(y)\dd y
        \quad(R\in\mathcal R).}
\end{equation}
This is the bridge between rational algebra and real integration.
First we check that the integrals make sense.

Put $f(y)=(e^{2\pi y}-1)^{-1}$. Its geometric series may be differentiated
term by term on $y\geq\varepsilon>0$, where the derivative series converge
uniformly. Four differentiations show that
\[
 w(y)=\frac{y^5}{12}f^{(4)}(y).
\]
Near zero, $f(y)=1/(2\pi y)+g(y)$ with $g$ analytic at zero. Consequently
$f^{(4)}(y)=24/(2\pi y^5)+O(1)$ and $w(y)\to1/\pi$.
For $y\geq1$, factor $e^{-2\pi y}$ out of the series in \eqref{eq:weight};
the remaining sum is bounded by
$\sum_{\ell\geq1}\ell^4e^{-2\pi(\ell-1)}<\infty$.
Thus $w(y)=O(y^5e^{-2\pi y})$ at infinity.

An admissible $R(y^2)$ is bounded near zero, because its poles lie away from
$[0,\infty)$. It grows at most polynomially as $y\to\infty$.
These two bounds on $w$ prove absolute convergence in
\eqref{eq:representation}.

For $R(t)=t^e$, the summands of the integral are nonnegative. Tonelli's
theorem allows us to interchange the sum and integral:
\begin{align*}
 \int_0^\infty y^{2e}w(y)\dd y
 &=\frac{(2\pi)^4}{12}\sum_{\ell\geq1}\ell^4
      \int_0^\infty y^{2e+5}e^{-2\pi\ell y}\dd y\\
 &=\frac{(2e+5)!}{12(2\pi)^{2e+2}}
      \sum_{\ell\geq1}\frac1{\ell^{2e+2}}.
\end{align*}
Here $\int_0^\infty y^m e^{-by}\dd y=m!/b^{m+1}$ for $b>0$:
substitution reduces this to $b=1$, and repeated integration by parts gives
$m!$. Euler's formula, proved in Appendix A, is
\[
 \zeta(2e+2)=
 \frac{(-1)^eB_{2e+2}(2\pi)^{2e+2}}{2(2e+2)!}.
\]
Substituting it above yields
\[
 \int_0^\infty y^{2e}w(y)\dd y
 =\frac{(-1)^eB_{2e+2}}{24}\,
       \frac{(2e+5)!}{(2e+2)!}
 =\mu_X(t^e).
\]
The factorial quotient is exactly $(2e+3)(2e+4)(2e+5)$.
The equality is independent of $X$, as the definition requires.

\subsection{Why the pole values are the same integrals}
\source{The Stieltjes representation is Pr\'evost's \cite[Theorem 1]{prevost}.
The fifth-order calculation is also in Fauzan \cite[Proposition 2.2]{fauzan}
and Lai \cite[Lemma 4.6]{lai}.}
The simple fractions need a different calculation. We do it for a real
number $a>0$, then specialize to a positive integer. Define
\[
 q_a(y)=\frac{y^5}{y^2+a^2}
       =y^3-a^2y+a^4\frac{y}{y^2+a^2}.
\]
The polynomial terms disappear after four derivatives. Since
$y/(y^2+a^2)=\operatorname{Im}(a-iy)^{-1}$, and
$\frac{d^4}{dy^4}(a-iy)^{-1}=24(a-iy)^{-5}$, we obtain
\[
 q_a^{(4)}(y)=24a^4\operatorname{Im}(a-iy)^{-5}.
\]
Using $w=y^5f^{(4)}/12$, four integrations by parts give
\begin{align*}
 \int_0^\infty\frac{w(y)}{y^2+a^2}\dd y
 &=\frac1{12}\int_0^\infty q_a(y)f^{(4)}(y)\dd y\\
 &=\frac1{12}\int_0^\infty q_a^{(4)}(y)f(y)\dd y\\
 &=2a^4\int_0^\infty
     \frac{\operatorname{Im}(a-iy)^{-5}}{e^{2\pi y}-1}\dd y.
\end{align*}
There are no missing boundary terms. First work on $[\varepsilon,L]$;
the boundary expression is
$q_af'''-q_a'f''+q_a''f'-q_a'''f$.
For $0\leq r\leq3$, at zero we have
\[
 q_a^{(r)}(y)=O(y^{5-r}),\qquad
 f^{(3-r)}(y)=O(y^{r-4}),
\]
so every boundary product is $O(y)$ and tends to zero. At infinity the
derivatives of $f$ decay exponentially and the other factors have at most
polynomial growth. Also $q_a^{(4)}(y)=O(y)$ at zero, so
$q_a^{(4)}(y)f(y)=O(1)$ there. The final integral converges at both ends,
and we may let $\varepsilon\downarrow0$ and $L\to\infty$.

For $a>0$, define $\zeta(5,a)=\sum_{v=0}^\infty(a+v)^{-5}$.
The particular Hurwitz identity needed here is
\begin{equation}\label{eq:hermite}
 \zeta(5,a)=\frac1{2a^5}+\frac1{4a^4}
 +2\int_0^\infty
     \frac{\operatorname{Im}(a-iy)^{-5}}{e^{2\pi y}-1}\dd y.
\end{equation}
Appendix A derives this directly from an elementary kernel identity.
It gives
\[
 \int_0^\infty\frac{w(y)}{y^2+a^2}\dd y
 =a^4\zeta(5,a)-\frac1{2a}-\frac14.
\]
For an integer $j\geq1$, the series gives
$\zeta(5,j)=\xi-\Hfive{j}+j^{-5}$. Therefore
\[
 \int_0^\infty\frac{w(y)}{y^2+j^2}\dd y
 =j^4(\xi-\Hfive{j})-\frac14+\frac1{2j}
 =\mu_\xi\!\left(\frac1{t+j^2}\right).
\]
We have checked both types of term in \eqref{eq:partial}.
Applying the formulas to its finitely many terms and using linearity
of integration proves \eqref{eq:representation} for every $R\in\mathcal R$.

\subsection{One transform, two kinds of data}
\source{This perspective is guided by Persiflage \cite{persiflage}.
The underlying transform and orthogonal-polynomial identities are
Pr\'evost's \cite[Section 2 and Theorems 1--2]{prevost}.}

It is helpful to see where the zeta value enters.
Change variables \(t=y^2\) and write
\[
 d\nu(t)=\frac{w(\sqrt t)}{2\sqrt t}\,dt,\qquad
 \Phi(z)=\int_0^\infty\frac{d\nu(t)}{z+t}\quad(z>0).
\]
The moments \(m_j=\int t^j\,d\nu(t)=\mu_X(t^j)\) are rational.
On the other hand, the pole formula gives
\[
 \Phi(a^2)=a^4\bigl(\xi-H_a^{(5)}\bigr)-\frac14+\frac1{2a}
 \qquad(a\in\Z_{>0}).
\]
Thus rational approximation to a transform value gives rational
approximation to \(\xi\).

The moments also describe \(\Phi\) for large \(z\).
The finite geometric identity gives
\[
 \Phi(z)=\sum_{j=0}^{m-1}\frac{(-1)^jm_j}{z^{j+1}}
  +\frac{(-1)^m}{z^m}\int_0^\infty\frac{t^m}{z+t}\,d\nu(t).
\]
The last term has absolute value at most \(m_m z^{-m-1}\).
This is an asymptotic expansion for each fixed \(m\); we have not asserted
convergence of the infinite moment series.

At \(X=\xi\), the two clauses defining \(\mu_X\) describe samples and
asymptotic coefficients of one transform. This is the bridge to
Pr\'evost's Pad\'e approximation, emphasized by Persiflage.
Our proof uses a rationally modified measure and its determinant;
no quantitative Pad\'e estimate is needed. The question is how much
real smallness survives when the determinant's coefficients are made
integral.

\section[Modifying the measure and its determinant]{Modifying the measure and taking its determinant}\label{sec:determinant}
\source{Fauzan \cite[Section 2.1]{fauzan}; Persiflage
\cite[Hankel Determinants]{persiflage}; Lai \cite[Section 5]{lai}.}
\subsection{The modifier and its affine moments}
For a positive integer $n$, choose
\[
 K=40n,\qquad N=3n,\qquad h=K-N=37n,
 \qquad D_m(t)=\prod_{j=1}^m(t+j^2).
\]
The parameters have been chosen to make the later estimates work.
For now, the features that matter are positivity and cancellation. Define
\[
 A_K(t)=\frac{D_N(t)^6}{D_K(t)}
       =\frac{D_N(t)^5}{T_K(t)},\qquad
 T_K(t)=\prod_{k=N+1}^K(t+k^2).
\]
The cancellation removes the first $N$ possible poles. The $h$ remaining
roots of $T_K$ are distinct negative integer squares, so multiplying
$A_K$ by any polynomial leaves a function in $\mathcal R$.
Moreover, $A_K(y^2)>0$ for every $y>0$.

It is useful to name the pairing before choosing its coordinates:
\begin{equation}\label{eq:bilinear-form}
 \mathcal B_{K,X}(p,q)=\mu_X(A_Kpq),
 \qquad p,q\in\Q[t]_{<h}.
\end{equation}
The monomials below and the later bases adapted to each prime are
coordinates for this one form. The real estimate uses its integral
representation. This common starting point is explicit in
Fauzan \cite[Sections 2.3--2.4]{fauzan}.

For $0\leq i,j<h$ put $R_{ij}(t)=A_K(t)t^{i+j}$ and define
\begin{equation}\label{eq:G}
 G_K(X)=\bigl[\mu_X(R_{ij})\bigr]_{0\leq i,j<h},
 \qquad \Delta_K(X)=\det G_K(X).
\end{equation}
The dependence on $i+j$ makes this a Hankel matrix: its entries are constant
along each anti-diagonal.
To see explicitly why these are affine entries, take partial fractions:
\[
 R_{ij}(t)=P_{ij}(t)+\sum_{k=N+1}^K\frac{c_k^{ij}}{t+k^2},
 \qquad
 c_k^{ij}=\frac{D_N(-k^2)^5(-k^2)^{i+j}}{T_K'(-k^2)}.
\]
The polynomial part may be zero; it must not otherwise be discarded.
Applying the functional gives
\[
 G_K(X)_{ij}=\mu_X(P_{ij})+
 \sum_{k=N+1}^K c_k^{ij}
 \left[k^4(X-\Hfive{k})-\frac14+\frac1{2k}\right].
\]
Thus every entry is $a_{ij}X+b_{ij}$ with rational $a_{ij},b_{ij}$.
Each term of the determinant expansion is a product of \(h\) affine
entries. Hence $\Delta_K\in\Q[X]$ and $\deg\Delta_K\leq37n$.
Section~\ref{sec:exact-degree} proves exact degree $h$, but the upper bound suffices.

\subsection{Evaluating the matrix at \texorpdfstring{\(\xi\)}{xi}}
\source{The positive Gram interpretation is explicit in Fauzan
\cite[Proposition 2.2]{fauzan}, Persiflage \cite{persiflage},
and Lai \cite[Proposition 5.6]{lai}.}
We can now pass from a symbolic matrix to a matrix of integrals one entry
at a time. With $R_{ij}$ as above,
\begin{align*}
 G_K(\xi)_{ij}
 &=\left.\mu_X(R_{ij})\right|_{X=\xi}\\
 &=\mu_\xi(R_{ij})\\
 &=\int_0^\infty R_{ij}(y^2)w(y)\dd y\\
 &=\int_0^\infty y^{2i}y^{2j}A_K(y^2)w(y)\dd y.
\end{align*}
The third equality is the representation proved in Section~\ref{sec:functional}.
The last is substitution: $(y^2)^{i+j}=y^{2i}y^{2j}$.
Absolute convergence was already checked for every admissible rational
function, including these particular $R_{ij}$.

For a compact description, set
\[
 \dd\nu_K(y)=A_K(y^2)w(y)\dd y,\qquad
 v(y)=(1,y^2,\ldots,y^{2h-2})^{\mathsf T}.
\]
Then $G_K(\xi)=\int_0^\infty v(y)v(y)^{\mathsf T}\dd\nu_K(y)$,
where the integral is taken \emph{entry by entry}. Its $(i,j)$ entry is
exactly the last integral above.

Let $c=(c_0,\ldots,c_{h-1})^{\mathsf T}\in\R^h$ and
$q(t)=\sum_i c_it^i$. Distributing the finite sums gives
\[
 c^{\mathsf T}G_K(\xi)c
 =\int_0^\infty\left(\sum_i c_iy^{2i}\right)^2\dd\nu_K(y)
 =\int_0^\infty q(y^2)^2\dd\nu_K(y).
\]
Only finite linearity is used here. Real coefficients $c_i$ cause no
difficulty with the originally $\Q$-linear definition of $\mu_X$.

If $c\ne0$, then $q$ is a nonzero polynomial. Choose $y_0>0$ that is not
a root of $q(y^2)$. By continuity, $q(y^2)^2$ is bounded below by a positive
number on a small interval about $y_0$. The density of $\nu_K$ is positive
and continuous there, so the integral is strictly positive.

Thus $G_K(\xi)$ is positive definite. To see why its determinant is positive,
orthogonalize $1,y^2,\ldots,y^{2h-2}$ by Gram--Schmidt without normalizing.
The resulting orthogonal functions are monic polynomials in $y^2$, so each
is nonzero and has positive squared norm. The change of basis is triangular
with diagonal entries one. The determinant is therefore the product of
these positive squared norms: $\Delta_K(\xi)>0$.

The matrix integral is not a rule for moving a determinant inside an
integral. Each $v(y)v(y)^{\mathsf T}$ has rank one. The multiple integral
for the determinant comes from Andr\'eief's identity, proved in the analytic section.

The integral and the affine dependence can now be displayed together.
For rational polynomials \(p,q\) of degree less than \(h\), put
\[
 d_k=\frac{k^4D_N(-k^2)^5}{T_K'(-k^2)}
 \qquad(N<k\le K).
\]
The partial fractions already computed give the exact identity
\begin{equation}\label{eq:gram-residue}
 \mathcal B_{K,X}(p,q)
 =\int_0^\infty p(y^2)q(y^2)\dd\nu_K(y)
 +(X-\xi)\sum_{k=N+1}^K d_kp(-k^2)q(-k^2).
\end{equation}
Indeed the polynomial part is independent of \(X\), while the change
in the value at the pole \(-k^2\) is \(k^4(X-\xi)\).
The residue sum is the slope; polynomial quotients still contribute to
the constant part. This explains the later arithmetic correction.
The identity combines Fauzan's Gram and residue formulas
\cite[Sections 2.3--2.4]{fauzan}.

\subsection{Putting both estimates on the same scale}
Fauzan uses the positive rational scalar
\begin{equation}\label{eq:normalization}
 S_K=\frac{(K!)^{2h}4^{h-1}}
 {(N!)^{12h}\prod_{i=1}^{h-1}((2i)!)^2},
 \qquad F_K(X)=S_K\Delta_K(X).
\end{equation}
This is his equation (2.5) \cite{fauzan}, also retained in
\texttt{Apery/Construction.lean} \cite{firsching}.
It follows at once that
\[
 F_K\in\Q[X],\qquad \deg F_K\leq h,\qquad F_K(\xi)>0.
\]
The rescaling is useful because factorial valuations admit exact formulas,
and factorial sizes admit Stirling estimates. It does not, on its own,
make every coefficient integral.

The goal itself is independent of rational scale. Let \(c(A)>0\)
be the rational content of a nonzero \(A\in\Q[X]\): dividing by \(c(A)\)
makes its coefficients integers with greatest common divisor one.
If \(v_p^{\mathrm G}(A)\) denotes the minimum coefficient valuation,
then \(c(A)=\prod_p p^{v_p^{\mathrm G}(A)}\), by taking a common
denominator and then the greatest common divisor. For \(A(\xi)\ne0\), set
\begin{equation}\label{eq:primitive-size}
 \mathcal J_\xi(A)
 =\log|A(\xi)|-\sum_p v_p^{\mathrm G}(A)\log p
 =\log\left|\frac{A(\xi)}{c(A)}\right|.
\end{equation}
For \(r\in\Q^\times\), the identity
\(\log|r|=\sum_p v_p(r)\log p\) proves
\(\mathcal J_\xi(rA)=\mathcal J_\xi(A)\).
The real change in size caused by a rational rescaling is exactly
balanced by its change at the finite primes.
Fauzan already uses primitive normalization
\cite[Appendix C.3]{fauzan}. We choose \(S_K\) for convenient estimates,
then construct an explicit multiplier from certified prime bounds.
Together those estimates bound \(\mathcal J_\xi\).

\section{Why the degree is exactly \texorpdfstring{\(h\)}{h}}\label{sec:exact-degree}
\source{Fauzan \cite[(2.9)]{fauzan}; Lai \cite[Proposition 6.1]{lai};
the same formula is stated in Persiflage \cite{persiflage}.}

The main proof only needs \(\deg\Delta_K\le h\). The residue pairing
also shows why the highest term survives: there are exactly \(h\)
remaining poles, matching the dimension of the polynomial space.
Put \(a_k=-k^2\) and \(V_{ik}=a_k^i\), with \(0\le i<h\) and \(N<k\le K\).
By \eqref{eq:gram-residue}, if \(G_K(X)=G_0+XG_1\), then
\[
 G_1=V\operatorname{diag}(d_k)V^{\mathsf T},\qquad
 d_k=\frac{k^4D_N(a_k)^5}{T_K'(a_k)}.
\]
Every \(d_k\) is nonzero. The nodes \(a_k\) are distinct, so evaluation
at them is injective on polynomials of degree less than \(h\):
a nonzero such polynomial cannot have \(h\) roots. Thus \(V\) is
invertible and \(\det G_1\ne0\). Taking the \(X\) term from every factor
in the determinant expansion gives
\[
 [X^h]\Delta_K(X)=\det G_1\ne0,\qquad \deg\Delta_K=h.
\]
For the explicit leading coefficient, the Vandermonde formula and
pairing the factors in the derivatives give
\[
 \det V=\prod_{k<\ell}(a_\ell-a_k),\qquad
 \prod_k T_K'(a_k)=(-1)^{h(h-1)/2}(\det V)^2.
\]
Hence
\[
 [X^h]\Delta_K(X)
 =(-1)^{h(h-1)/2}\prod_{k=N+1}^K k^4D_N(-k^2)^5
 \in\Z\setminus\{0\}.
\]
This concerns one coefficient; the arithmetic must still control them all.

\section{The analytic estimate: one bound for every configuration}
\label{analytic:section}

The real estimate has three steps: express the determinant as a positive
integral, bound its logarithm by an energy comparison, and restore the
factorial normalization. Fauzan supplies the method and comparison
measure \cite[Section 6 and Appendix A]{fauzan}; Firsching's analytic
files give a further implementation \cite{firsching}.
The detailed companion, Sections 5.1--5.6, proves the supporting bounds
and records the finite certificate. Here we keep their role in the
argument visible.
Throughout, \(\alpha=N/K=3/40\) and \(\lambda=h/K=37/40\).

\subsection{From the determinant to logarithmic energy}

Applying Andr\'eief's identity to the positive Gram matrix gives
\begin{equation}\label{analytic:andreief}
 \Delta_K(\xi)=\frac1{h!}\int_{(0,\infty)^h}
 \prod_{i<j}(y_i^2-y_j^2)^2
 \prod_i A_K(y_i^2)w(y_i)\,dy_1\cdots dy_h.
\end{equation}
Expanding the two Vandermonde determinants and integrating term by term
proves the identity: relabeling the variables leaves \(h!\) copies of
the Gram determinant. Exponential decay makes every term integrable
\cite{andreief}.

Put \(y_i=K\sqrt{t_i}\). The external field comes directly from the
finite products:
\[
 \log D_m(K^2t)
 =2m\log K+K\int_0^{m/K}\log(t+u^2)\,du+O(\log K),
\]
uniformly for \(t\ge0\) and \(1\le m\le K\).
The error is between zero and \(2\log K+2\), by comparison with the
right-endpoint Riemann sum (companion, Section 5.5).
Including the exponential factor in \(w\) gives
\begin{equation}\label{analytic:field}
 V(t)=2\pi\sqrt t+\int_0^1\log(t+u^2)\,du
                 -6\int_0^\alpha\log(t+u^2)\,du.
\end{equation}
Thus the main exponent of the integrand is
\[
 E_K(\mathbf t)=2\sum_{i<j}\log|t_i-t_j|-K\sum_i V(t_i).
\]
The first term rewards separation; the second confines the points.
This is the real counterpart of taking valuations of the same products
in the arithmetic estimate.

\subsection{Why one comparison measure is enough}

Write \(U^\rho(t)=\int\log|t-u|\,d\rho(u)\) and
\(I(\rho)=\iint\log|t-u|\,d\rho(t)d\rho(u)\).
Fauzan's explicit positive combination of sixteen arcsine measures
has mass \(\lambda\), support in \((0,2)\), and satisfies
\begin{equation}\label{analytic:potential-target}
 2U^\rho(t)-V(t)\le M_0=-\frac{1329}{200}
 \qquad(t\ge0).
\end{equation}
Its potential is certified on all 684 intervals of a partition of
\([0,2]\), using outward rational bounds; a decreasing elementary
majorant handles \(t\ge2\). The data and verification are in the
companion, Sections 5.2--5.3, and \texttt{checks/analytic\_certificate.py}.
This measure is chosen to prove a useful inequality, not to solve
the exact equilibrium problem.

The reason the inequality works is short. With this \(\rho\), let
\(\eta\) be positive and compactly supported on \([0,\infty)\), of
mass \(\lambda\), with absolute logarithmic integrability for
\((\eta+\rho)\otimes(\eta+\rho)\).
The zero-mass energy inequality gives \(I(\eta-\rho)\le0\), hence
\[
 I(\eta)-\int V\,d\eta
 =I(\eta-\rho)+\int(2U^\rho-V)\,d\eta-I(\rho)
 \le\lambda M_0-I(\rho).
\]
The Gaussian-kernel proof of the zero-mass inequality is in the
companion, Section 5.4. The original moment measure defines the
determinant; \(\eta\) describes a configuration; \(\rho\) bounds all
configurations at once.

The configuration measure \(\eta=K^{-1}\sum_i\delta_{t_i}\) has mass
\(\lambda\). With the diagonal omitted,
\(I(\eta)-\int V\,d\eta=E_K/K^2\).
Its self-energy is infinite and negative, so it cannot be inserted
directly. Replace each point mass by arclength
on a circle of radius \(K^{-2}\). The circle self-energies and the
change in the real-axis potential produce only \(O(K\log K)\).
Keeping an extra decaying factor gives the precise usable bound
\begin{align}\label{analytic:configuration}
 E_K(\mathbf t)+\sum_i\sqrt{t_i}
 \le{}&(\lambda M_0-I(\rho))K^2\\
      &+(120+\sqrt2)h+2h\log K,\notag
\end{align}
for every \(t_i\ge0\) and admissible \(K\).
The companion's Section 5.4 proves the circle-average and tail
estimates; a real-axis potential inequality alone would not suffice.
The retained factor \(\exp(-\sum_i\sqrt{t_i})\) makes the bound integrable.

\subsection{The normalization cancels the larger terms}

The remaining factors are elementary. The bound
\(w(y)\le8192(1+y)^5e^{-2\pi y}\), the Jacobians, and the product
errors contribute \(O(K\log K)\) to the logarithm. After
\eqref{analytic:configuration}, the remaining one-variable integral is
\(\int_0^\infty t^{-1/2}(1+\sqrt t)^5e^{-\sqrt t}\,dt=652\).
The powers of \(K\) in the Vandermonde and \(A_K\) give
\(\kappa=2\lambda(\lambda+6\alpha-1)\), so
\[
 \log\Delta_K(\xi)
 \le\kappa K^2\log K+(\lambda M_0-I(\rho))K^2+O(K\log K).
\]
The factorial bounds for \(S_K\) give the opposite leading term:
\begin{align*}
 \log S_K&\le-\kappa K^2\log K+C_*K^2+O(K\log K),\\
 C_*&=-2\lambda+12\alpha\lambda(1-\log\alpha)
             +3\lambda^2-2\lambda^2\log(2\lambda).
\end{align*}
Thus normalization cancels the \(K^2\log K\) terms exactly. The
certificate verifies
\[
 \lambda M_0-I(\rho)+C_*<-1.3669955.
\]
Keeping the elementary errors explicit yields, for every \(K=40n\),
\begin{equation}\label{analytic:goal}
 \boxed{\ \log F_K(\xi)
 \le-1.3669955K^2+27K\log K+200K.\ }
\end{equation}
Sections 5.5--5.6 of the companion carry out this bookkeeping.
The important margin is the negative quadratic coefficient; the
arithmetic must make the coefficients integral at a smaller
quadratic cost.

\section{Arithmetic: choosing coordinates at each prime}
\label{sec:local-arithmetic}

The arithmetic must control every coefficient of \(F_K=S_K\Delta_K\).
For \(A(X)=\sum_j a_jX^j\), write
\[
 v_p^{\mathrm G}(A)=\min_{a_j\ne0}v_p(a_j),\qquad
 v_p^{\mathrm G}(0)=+\infty.
\]
Thus \(v_p^{\mathrm G}(A)\ge L\) means that \(p^L\) divides all its
coefficients, with negative \(L\) allowing denominators. Evaluating
at a particular \(X\) would not give this information.
The strategy is Fauzan's; the bounds used here are Firsching's revised
finite partial-fraction estimates
\cite[Sections 3--4]{fauzan}\cite[\texttt{Arith/}, \texttt{Norm.lean}]{firsching}.

Two elementary facts organize the calculation. A basis change with
coefficient matrix \(C\) replaces \(G_K\) by \(C^{\mathsf T}G_KC\);
its determinant is multiplied by \((\det C)^2\).
Hence a basis whose determinant is a \(p\)-adic unit preserves the
valuation. Also,
\begin{equation}\label{arith:detweights}
 v_p^{\mathrm G}(B_{ij})\ge w_i+w_j
 \quad\Longrightarrow\quad
 v_p^{\mathrm G}(\det B)\ge2\sum_iw_i.
\end{equation}
Every product in the determinant expansion has that lower bound.
The task is therefore to choose coordinates in which useful row
weights are visible.

\subsection{Three bases for the same form}

All three constructions estimate
\(\mathcal B_{K,X}(f,g)=\mu_X(A_Kfg)\).
Their difference is the arithmetic information they expose.

\emph{Small primes: integer values.}
The identity
\[
 \mu_X(R)=\tau_X\bigl(x^5R(-x^2)\bigr),\qquad
 \tau(P)=\frac1{24}\mathcal L(P'''),\qquad \mathcal L(x^j)=B_j,
\]
puts the poles at integers and expresses the polynomial contribution
through the same Bernoulli data as the real moments
\cite[(3.1)]{fauzan}. Appendix~\ref{app:pullback} defines \(\tau_X\)
and proves the identity.
In a binomial basis, integer values give integral finite differences.
The identities for differentiating \(\binom{x}{j}\) and
\(\mathcal L(\binom{x}{j})=(-1)^j/(j+1)\) then bound the loss of
divisibility. Applied to the specially chosen integer-valued basis,
whose determinant accounts for \(S_K\), this gives the uniform bound
\begin{equation}\label{arith:small}
 v_p^{\mathrm G}(F_K)\ge
 -6h\lfloor\log_p(5K)\rfloor-hv_p(24).
\end{equation}
The proof, including the polynomial parts of partial fractions, is
in the detailed companion, Sections 6.2--6.3.

\emph{Inner primes: distribute vanishing among residue classes.}
For \(K/400<p\le K/3\), several poles can occupy one class.
Choose polynomials with prescribed orders of vanishing at
\(t=-c^2\pmod p\). The Chinese remainder theorem makes them a basis
with unit determinant. Allocating the \(h\) basis vectors among these
classes balances zeros against poles.

Each row's weight must be valid in \emph{every} class. This is why the
raw vanishing-order weights are capped at a common level; the zero
class also needs a small reserve of rows because \(t=-x^2\) doubles
its vanishing orders. If \(W_p^{\rm in}\) is the sum of the resulting
capped weights, the entry bounds and \eqref{arith:detweights} give
\[
 v_p^{\mathrm G}(\Delta_K)\ge2W_p^{\rm in}.
\]
The explicit allocation and its uniform error bound occupy companion
Sections 6.4--6.6. They implement Fauzan's class basis with Firsching's
capped weights \cite[\texttt{InnerWeights.lean},
\texttt{InnerAlloc.lean}]{firsching}.

\emph{Outer primes: cancel poles, then restore the polynomial part.}
For \(K/3<p\le2h\), a basis built from the actual pole factors cancels
most troublesome denominators. The remaining low-order part has
nonpositive row weights with sum \(W_p^{\rm out}\).
The polynomial terms not covered by this estimate form a correction
of rank at most
\[
 r_p=\max(0,52n+2-2p).
\]
In the chosen basis the full matrix has the form \(A+p^{-1}UV\),
where \(U,V\) are integral and factor through \(r_p\) coordinates.
Indeed,
\[
 \det(A+p^{-1}UV)=
 \det\begin{pmatrix}A&-U\\p^{-1}V&I_{r_p}\end{pmatrix}.
\]
Applying the row-and-column version of \eqref{arith:detweights}
to this block matrix gives
\begin{equation}\label{arith:outer}
 v_p^{\mathrm G}(\Delta_K)\ge2W_p^{\rm out}-r_p.
\end{equation}
This pays the \emph{full} rank loss. The pole cancellations and the rank
calculation are in companion Section 6.7
\cite[\texttt{OuterFrame.lean}, \texttt{LowRank.lean}]{firsching}.

These are coefficient bounds for one determinant, in three coordinate
systems. The inner and outer constructions apply uniformly once
\(K>480000\); a finite initial range is irrelevant to the asymptotic
contradiction. The remaining task is to add their costs over primes.

\section{Adding the prime costs}\label{growth:section}

Fix \(M=400\), \(\alpha=3/40\), and \(\lambda=37/40\).
For sufficiently large \(K=40n\), the preceding bounds give integers
\begin{equation}\label{growth:Lp}
 L_p=
 \begin{cases}
 -6h\lfloor\log_p(5K)\rfloor-hv_p(24),&p\le K/400,\\
 v_p(S_K)+\lfloor2W_p^{\rm in}\rfloor,&K/400<p\le K/3,\\
 v_p(S_K)+\lfloor2W_p^{\rm out}-r_p\rfloor,&K/3<p\le2h,
 \end{cases}
\end{equation}
with \(L_p\le v_p^{\mathrm G}(F_K)\).
For \(p>2h\), the same outer calculation has no large-pole rows or
rank loss, and the factorial scalar is a unit; hence
\(v_p^{\mathrm G}(F_K)\ge0\).
This is the normalizer in Firsching's revised argument
\cite[\texttt{Norm.lean}]{firsching}.

Set
\begin{equation}\label{growth:normalizer}
 m_K=\prod_{p\le2h}p^{-L_p}.
\end{equation}
Every coefficient of \(m_KF_K\) has nonnegative valuation at every
prime, so its reduced denominator is one. Consequently
\begin{equation}\label{growth:integrality}
 m_KF_K\in\Z[X],\qquad
 \log m_K=\sum_{p\le2h}(-L_p)\log p.
\end{equation}
The multiplier is positive and rational. Positive \(L_p\)'s allow
common factors to be divided out and contribute \emph{negative} costs;
discarding those savings would spoil the estimate.

\subsection{From local counts to a fixed budget}

The local formulas still depend on \(K\) and \(p\).
Writing \(x=K/p\) converts them into bounded, piecewise continuous
profiles on fixed intervals. First the residue classes \(c/p\)
are averaged; then the resulting costs are summed over primes.
The profiles retain both the basis weights and the factorial
valuation \(v_p(S_K)\), computed by Legendre's formula.
Their explicit definitions are in companion Sections 7.4 and 7.7.
For each range above the small-prime cutoff, the corresponding profile
\(E\) satisfies
\[
 -L_p\le pE(K/p)+O_{400}(1).
\]
This converts the local divisibility bound into a cost that can be summed.

The prime number theorem,
\(\vartheta(T)=\sum_{p\le T}\log p=T+o(T)\), gives
\begin{equation}\label{growth:prime-sum}
 \frac1{K^2}\sum_{K/b<p\le K/a}p\log p\,E(K/p)
 \longrightarrow\int_a^b\frac{E(x)}{x^3}\,dx
\end{equation}
for a bounded piecewise continuous \(E\) with finitely many
discontinuities on a fixed \(0<a<b\).
To see the factor \(x^{-3}\), the prime number theorem replaces
\(K^{-1}\sum(\log p)\delta_{p/K}\) by length measure; apply this to
\(tE(1/t)\), then substitute \(x=1/t\).
An \(O_{400}(1)\) error at each prime contributes only
\(O_{400}(\vartheta(2h))=O(K)=o(K^2)\).
This is the prime-summation step of Fauzan and Firsching
\cite[Section 5]{fauzan}\cite[\texttt{PrimeSumIntegral.lean}]{firsching}.

There are four contributions. Small primes use the coarse bound:
apart from \(O(\sqrt K)\) primes, the logarithmic floor equals one,
giving \(6\lambda/400=0.013875\).
For \(3\le x\le20\), the inner profile has 125 rational pieces,
integrated exactly. For \(20\le x\le400\), a quadratic bound followed
by two integrations by parts gives a negative contribution.
The successive primitives are periodic Bernoulli polynomials: they
remove the oscillating fractional-part counting error.
The outer profile has ten merged rational pieces and includes the
full rank correction. Companion Sections 7.3--7.8 give these calculations;
the supplied exact-arithmetic checks reproduce their values.

\begin{center}
\begin{tabular}{l r}
\toprule
Prime range & Upper bound for its \(K^2\) coefficient\\
\midrule
\(p\le K/400\) & \(0.013875\)\\
\(K/400<p\le K/20\) & \(-0.05346\)\\
\(K/20<p\le K/3\) & \(0.04279\)\\
\(K/3<p\le2h\) & \(1.34481\)\\
\bottomrule
\end{tabular}
\end{center}
These finite decimals are rational upper bounds, rounded upwards
from the exact values. Their sum is \(1.348015<1.35\).
The exact calculation gives the slightly stronger bound
\[
 \limsup_{\substack{K\to\infty\\40\mid K}}
 \frac{\log m_K}{K^2}
 \le B=1.348003127429204\ldots<1.35.
\]
The exact fraction defining \(B\) is recorded with the final comparison
below and in the companion's Section 7.8.
Because the gap below \(1.35\) is strict, all sufficiently large
admissible \(K\) satisfy
\begin{equation}\label{growth:final}
 \boxed{m_K>0,\qquad m_KF_K\in\Z[X],\qquad
        \log m_K\le1.35K^2.}
\end{equation}
This is the arithmetic budget to compare with the real decay.

\section{Combining the estimates}\label{sec:conclusion}

Every estimate has been made for the same rational polynomial
\(F_K=S_K\Delta_K\). The analytic bound gives
\[
 \log F_K(\xi)\leq -1.3669955K^2+O(K\log K).
\]
Since \(K\log K=o(K^2)\), this implies
\[
 0<F_K(\xi)\leq \exp(-1.36K^2)
\]
for all sufficiently large admissible \(K\).
The arithmetic construction gives \(m_K>0\), \(m_KF_K\in\Z[X]\),
and
\[
 \limsup_{\substack{K\to\infty\\40\mid K}}
       \frac{\log m_K}{K^2}
 \leq
 B=
 \frac{2031621487545722214389669}{1507134105408387456000000}
 =1.348003127429\ldots<1.35.
\]
The strict gap implies \(\log m_K\leq1.35K^2\) for all sufficiently
large \(K\). Indeed, take any positive
\(\varepsilon<1.35-B\) in the defining upper-limit estimate.

Set \(Q_n(X)=m_{40n}F_{40n}(X)\). Its coefficients are integers.
Both scalar factors are positive, and the determinant was proved positive
at \(\xi\); hence \(Q_n(\xi)>0\). Its degree is at most \(37n\)
(and in fact equals \(37n\)).
For all sufficiently large \(n\),
\[
 \log Q_n(\xi)
 =\log m_{40n}+\log F_{40n}(\xi)
 \leq (1.35-1.36)(40n)^2=-16n^2.
\]
Thus
\[
 Q_n\in\Z[X],\qquad
 \deg Q_n=37n,\qquad
 0<Q_n(\zeta(5))\leq e^{-16n^2}.
\]
The integer contradiction in Section~\ref{sec:orientation} proves
that \(\zeta(5)\) is irrational.

Equivalently, \(\mathcal J_\xi(\Delta_K)=\mathcal J_\xi(F_K)\le-0.01K^2\)
eventually. Rationality would force \(\mathcal J_\xi(F_K)\ge-h\log b\):
a fixed rational denominator can offset only linear decay.

The constant 16 is a convenient rounding of the estimates established
here. It is not the original paper's claimed rate, and it is not a claim
about the stronger \(K=12n\) construction mentioned in Persiflage.
The revised normalization, the comparison measure, and the finite
certificates have all been kept consistent throughout.

\appendix
\section{Classical formulas used in the integral representation}\label{app:classical}
\source{Euler's formula, Hermite's formula, and the partial fractions of
\(\coth\) are classical; exact references are \cite{dlmf}.
This appendix expands the derivation in the earlier draft of this chat.}

The moment and pole calculations use two familiar facts: Euler's formula
for even zeta values and a special Hurwitz zeta integral. Both follow from
one elementary identity. Put $f(y)=(e^{2\pi y}-1)^{-1}$. For $u>0$,
\begin{equation}\label{eq:kernel}
 \frac1{1-e^{-u}}=\frac12+\frac1u
       +2\int_0^\infty f(y)\sin(uy)\dd y.
\end{equation}

To prove it, periodically extend
$F_u(x)=\cosh(u(x-\tfrac12))$ from $[0,1]$. Its endpoints agree,
and its Fourier coefficients satisfy
\[
 c_n=\int_0^1 F_u(x)e^{-2\pi\mathrm i nx}\dd x
     =\frac{2u\sinh(u/2)}{u^2+4\pi^2n^2},\qquad n\in\Z.
\]
Indeed, integrate $F_u''=u^2F_u$ twice by parts. The equal endpoint
values cancel, and the endpoint derivative difference is
$2u\sinh(u/2)$, giving
$(u^2+4\pi^2n^2)c_n=2u\sinh(u/2)$.
The periodic function is continuous and piecewise smooth. The Fourier
convergence theorem at $x=0$ therefore gives
$\cosh(u/2)=c_0+2\sum_{n\geq1}c_n$. Dividing by $2\sinh(u/2)$ yields
\[
 \frac12\coth(u/2)=\frac1u
      +2u\sum_{n\geq1}\frac1{u^2+4\pi^2n^2}.
\]
Expanding $f$ as a geometric series and integrating each exponential gives
\[
 \int_0^\infty f(y)\sin(uy)\dd y
 =\sum_{n\geq1}\int_0^\infty e^{-2\pi ny}\sin(uy)\dd y
 =\sum_{n\geq1}\frac{u}{u^2+4\pi^2n^2}.
\]
This interchange is justified by $|\sin(uy)|\leq uy$:
the sum of the absolute integrals is at most
$u\sum_{n\geq1}(2\pi n)^{-2}<\infty$.
Together with $1/(1-e^{-u})=\tfrac12+\tfrac12\coth(u/2)$,
this proves \eqref{eq:kernel}.

\textbf{Euler's formula.}
The same partial fractions give, for $0<u<2\pi$,
\begin{align*}
 \frac{u}{e^u-1}
 &=1-\frac u2+2u^2\sum_{n\geq1}\frac1{u^2+4\pi^2n^2}\\
 &=1-\frac u2+
   2\sum_{r\geq0}\frac{(-1)^r\zeta(2r+2)}{(2\pi)^{2r+2}}u^{2r+2}.
\end{align*}
For the second equality, expand each denominator geometrically. The
resulting double series converges absolutely whenever $|u|<2\pi$, so
the order of summation may be exchanged. Compare with the defining
series $u/(e^u-1)=\sum_{k\geq0}B_ku^k/k!$. Equality of coefficients gives
\[
 \boxed{\zeta(2e+2)=
 \frac{(-1)^eB_{2e+2}(2\pi)^{2e+2}}{2(2e+2)!}},
 \qquad e\in\Z_{\geq0}.
\]
This is the identity used to turn the polynomial integrals into the
rational Bernoulli values defining $\mu$.

\subsection{The Hurwitz integral}

Fix $a>0$. The definition
$\zeta(5,a)=\sum_{n\geq0}(a+n)^{-5}$ and the elementary exponential
integral give
\[
 \frac1{(a+n)^5}=\frac1{24}
       \int_0^\infty u^4e^{-(a+n)u}\dd u.
\]
The integrands are nonnegative, so Tonelli's theorem allows the sum and
integral to be exchanged. Summing the geometric series gives
\[
 \zeta(5,a)=\frac1{24}
       \int_0^\infty\frac{u^4e^{-au}}{1-e^{-u}}\dd u.
\]
Insert \eqref{eq:kernel}. The first two terms can be evaluated immediately:
\[
 \frac1{24}\int_0^\infty u^4e^{-au}
       \left(\frac12+\frac1u\right)\dd u
       =\frac1{2a^5}+\frac1{4a^4}.
\]
The remaining term is a double integral. Although the sine can change
sign, Fubini's theorem applies because
\begin{align*}
 &\int_0^\infty\!\int_0^\infty
       u^4e^{-au}f(y)|\sin(uy)|\dd y\dd u\\
 &\hspace{12mm}\leq
       \left(\int_0^\infty u^5e^{-au}\dd u\right)
       \left(\int_0^\infty yf(y)\dd y\right)<\infty.
\end{align*}
The first factor is $5!/a^6$. For the second, $yf(y)$ tends to
$1/(2\pi)$ at zero and decays exponentially at infinity. Thus the bound
also justifies subtracting and rearranging the integrals in the preceding
calculation.

After exchanging the integrations, evaluate the inner integral by taking
the imaginary part of an exponential integral:
\[
 \int_0^\infty u^4e^{-au}\sin(uy)\dd u
 =\operatorname{Im}\int_0^\infty
       u^4e^{-(a-\mathrm i y)u}\dd u
 =24\operatorname{Im}(a-\mathrm i y)^{-5}.
\]
The last equality follows by four integrations by parts; every boundary
term vanishes because the real part of $a-\mathrm i y$ is positive.
Keeping the factor $2/24$ from \eqref{eq:kernel}, we obtain
\[
 \boxed{\zeta(5,a)=\frac1{2a^5}+\frac1{4a^4}
 +2\int_0^\infty
   \frac{\operatorname{Im}(a-\mathrm i y)^{-5}}{e^{2\pi y}-1}\dd y.}
\]
This is precisely the special Hermite formula \eqref{eq:hermite} used in
the pole calculation. In particular, it holds for every positive integer
$a=j$, where the Hurwitz zeta value is the tail of the series for
$\zeta(5)$.

The two classical identities have now been derived from the same kernel.
The remaining passage to the positive weight $w$ is the integration by
parts carried out in the main text. Together, these calculations justify
the integral representation on both parts of the partial-fraction basis.

\bigskip
{\small\textbf{Classical references.}
The identities derived here are also recorded in the NIST Digital Library
of Mathematical Functions:
\href{https://dlmf.nist.gov/4.36.E3}{the partial fractions for $\coth$ (4.36.3)},
\href{https://dlmf.nist.gov/25.6.E2}{Euler's even-zeta formula (25.6.2)}, and
\href{https://dlmf.nist.gov/25.11.E29}{Hermite's Hurwitz integral (25.11.29)}.}

\section{The pullback to integer poles}\label{app:pullback}
\source{Fauzan \cite[(3.1)]{fauzan}; Lai \cite[Lemma 7.3]{lai};
Firsching \cite[\texttt{Arith/Pullback.lean}]{firsching}.}

The auxiliary arithmetic functional changes the pole locations from
negative squares to integers. Here is the complete identity, including
the weight that accompanies the substitution.

Define $L:\Q[x]\to\Q$ by $L(x^k)=B_k$, and put
\[
 \tau(P)=\frac{L(P''')}{24},\qquad
 d(r)=\begin{cases}r,&r\geq0,\\-r-1,&r<0.\end{cases}
\]
On rational functions with only simple integer poles, define $\tau_X$
by its polynomial restriction $\tau$ and the pole values
\[
 \tau_X\!\left(\frac1{x-r}\right)=\Hfive{d(r)}-X,
 \qquad r\in\Z.
\]
Unique partial fractions extend these prescriptions linearly, just as
in the definition of $\mu_X$.

For every rational function $R$ in the domain of $\mu_X$, the identity is
\begin{equation}\label{eq:pullback}
 \boxed{\mu_X(R)=\tau_X\!\left(x^5R(-x^2)\right).}
\end{equation}
The right side is defined: a pole at $t=-j^2$ becomes a pair of simple
poles at $x=j$ and $x=-j$, and a polynomial remains a polynomial after
substitution and multiplication by $x^5$. Both sides are linear in $R$.
By \eqref{eq:partial}, it is enough to check monomials and simple poles.

For $R(t)=t^e$, the transformed function is $(-1)^ex^{2e+5}$. Three
differentiations give
\[
 \tau_X\bigl((-1)^ex^{2e+5}\bigr)
 =\frac{(-1)^e(2e+5)(2e+4)(2e+3)B_{2e+2}}{24},
\]
which is exactly \eqref{eq:mu-moment}. This explains the order of the
derivative: it produces both the three consecutive factors and the
required Bernoulli index.

For $R(t)=1/(t+j^2)$, polynomial division gives
\[
 \frac{x^5}{j^2-x^2}
 =-x^3-j^2x-\frac{j^4}{2}
       \left(\frac1{x-j}+\frac1{x+j}\right).
\]
One can check the signs by multiplying by $j^2-x^2$. The polynomial
terms give $x^5-j^4x$, and the pole terms give $+j^4x$.
Now $\tau(-x^3)=-1/4$, $\tau(-j^2x)=0$, $d(j)=j$, and $d(-j)=j-1$.
Applying $\tau_X$ therefore gives
\begin{align*}
 \tau_X\!\left(\frac{x^5}{j^2-x^2}\right)
 &=-\frac14-\frac{j^4}{2}
       \left(\Hfive{j}+\Hfive{j-1}-2X\right)\\
 &=j^4\bigl(X-\Hfive{j}\bigr)-\frac14+\frac1{2j},
\end{align*}
where the last equality uses
$\Hfive{j}-\Hfive{j-1}=j^{-5}$. This is \eqref{eq:mu-pole}, completing
the proof by linearity.

Both the substitution $t=-x^2$ and the factor $x^5$ are necessary.
The identity supplies a second description of the same functional;
it introduces no new determinant. Integer pole locations make the
arithmetic calculations by residue class possible, while the original
$t$-description makes the positive integral representation visible.

\section{Attribution, dependencies, and verification}\label{app:attribution}

This draft was prepared with OpenAI Codex at the user's request.
The earlier drafts in the same chat supplied some of the expanded
elementary calculations. Their underlying mathematical sources remain
Fauzan, Lai, Pr\'evost, and the classical identities listed below.
Reusing those explanations does not make their mathematics original.

This revision also draws on the chat
\emph{Find parallel constructions in the Zeta(5)\ldots}
and its corrected explanatory companion.
They suggested making the common bilinear form, the roles of the local
bases, the passage from finite products and counts to continuous
estimates, and the primitive evaluation size explicit.
The chat's final provenance correction is retained: Fauzan already
states the underlying connections, including the weighted pullback,
the local-basis strategy, the periodic primitives, and the comparison
measure argument. The chat contributes organization and classical
interpretations. It supplies no new quantitative estimate.

\begin{longtable}{>{\raggedright\arraybackslash}p{0.27\textwidth}>{\raggedright\arraybackslash}p{0.65\textwidth}}
\toprule
Part of this account & Source and nature of the dependence\\
\midrule
\endhead
Guiding narrative &
Calegari's Persiflage post: the positive-measure, rational-modifier,
Hankel-determinant viewpoint and the connection with prior approximation
theory. It precedes Lai's exposition.\\[3pt]
Positive measure and transform &
Pr\'evost, Theorems 1--2. Euler's and Hermite's classical identities
are also recorded in DLMF. The special fifth-order calculation is
expanded here, as in Fauzan and Lai.\\[3pt]
Concrete determinant &
Fauzan, Section 2: \(D_N^6/D_K\), \(K=40n,N=3n\), affine entries,
degree, positivity, and the factorial scalar. Lai is a cross-check
and an exposition used in the preceding drafts.\\[3pt]
General determinant context &
Zudilin's irrationality criterion; Brown's integral pairing and
normalization framework; Uvarov's rational modifications, as discussed
and extended by Krattenthaler. These are contextual antecedents,
not additional unproved estimate inputs.\\[3pt]
Analytic estimate &
Fauzan, Section 6 and Appendix A: external field, comparison measure,
configuration inequality, and interval partition. Firsching's analytic
files supply detailed cross-checks. The finite certificate was re-executed
for the detailed version; its checked files accompany this condensation.\\[3pt]
Local arithmetic &
Fauzan's weighted-basis strategy, in the revised finite
partial-fraction implementation of Firsching's \texttt{Arith} files.
The full rank loss and cutoff 400 are retained.\\[3pt]
Global arithmetic &
Firsching's \texttt{Norm} and \texttt{Growth} files: allocation,
prime ranges, continuous profiles and corrected numerical total.
The rational pieces and integrals were independently regenerated
for the detailed version; the same data are used here.\\[3pt]
Parallel-construction explanations &
The named chat and its revised companion: the common-pairing
presentation, Euler--Maclaurin language, periodic Bernoulli notation,
and scale-independent evaluation size. Their mathematical inputs
are credited to Fauzan, the classical Bernoulli identities, and
Firsching's revised estimates at the relevant steps.\\[3pt]
Finite computation &
The companion files identify which data are transcribed from Fauzan
or Firsching and which verification code was written for this draft.
Any adapted repository code retains its license and attribution.\\
\bottomrule
\end{longtable}

The local source ledger records URLs, versions, hashes, and precise
locations. The mathematical proof and its finite checks should be
distinguished from a complete local rebuild of the formal development:
no such rebuild is claimed. The source repository's CI results are
external evidence, not a substitute for the calculations described here.

The uniform potential bound is certified interval by interval, including
the range beyond the finite partition. The arithmetic tables are checked
with rational arithmetic on complete intervals. Neither certificate is
based on testing selected floating-point values.

The background theorems are Tonelli and Fubini, elementary Fourier
convergence, finite dimensional linear algebra, the prime number theorem,
and the classical Bernoulli and factorial facts. Argument-specific
supporting estimates cited in the concise sections are proved in the
separate detailed companion.

\begin{thebibliography}{99}
\addcontentsline{toc}{section}{References}

\bibitem{persiflage}
Frank Calegari (Persiflage),
\emph{\(\zeta(5)\) is irrational}, 24 September 2026.
\href{https://galoisrepresentations.org/2026/09/24/zeta5-is-irrational/}{Blog post}.
The guiding exposition for this draft. The displayed post date is earlier
than Lai's write-up; no dependence of the post on Lai is asserted.
Calegari identifies Persiflage as his blog on his
\href{https://math.uchicago.edu/~fcale/research.html}{research page}.

\bibitem{fauzan}
Aabir Fauzan,
\emph{\(\zeta(5)\) is irrational}, manuscript dated 17 September 2026,
32 pages, Zenodo record 22826419.
\href{https://zenodo.org/records/22826419}{DOI: 10.5281/zenodo.22826419}.
The determinant, factorial normalization, local arithmetic strategy,
comparison measure, and analytic certificate originate here.

\bibitem{firsching}
Moritz Firsching and contributors,
\emph{Zeta5: Lean formalization of the irrationality of \(\zeta(5)\)}.
\href{https://github.com/mo271/Zeta5/tree/7fe736760f4b96bfdb4334b68e3b3124ecbe10b0}{Repository at commit
\texttt{7fe7367}}, 27 September 2026. Apache License 2.0.
The substantive initial formalization is commit
\href{https://github.com/mo271/Zeta5/commit/2c80b320d06b4395605edbd413649d19094507af}{\texttt{2c80b32}}.
The revised local estimates, cutoff 400, full rank loss, and global
normalization used in this draft are translated from the
\texttt{Apery/Arith}, \texttt{Apery/Norm.lean}, and
\texttt{Apery/Growth} files, with exact locations in the source ledger.
The formal development depends on Lean, Mathlib and
\href{https://github.com/mo271/PrimeNumberTheoremAnd/tree/04d81b2a68ea3548b9990813fdc463c78511f412}{PrimeNumberTheoremAnd};
their contributors retain credit for those libraries and formalized
background theorems. Paul-Lez is recorded as a coauthor of the Comparator
CI addition. No local rebuild of the complete Lean development is claimed.

\bibitem{lai}
Chun Wing Lai,
\emph{On Fauzan's Proof that \(\zeta(5)\) is Irrational},
v1, 29 September 2026.
\href{https://ai.vixra.org/pdf/2609.0080v1.pdf}{ai.vixra:2609.0080v1}.
Used as a later exposition of the construction and as a cross-check;
its printed arithmetic conclusion is not an input to this draft.

\bibitem{prevost}
Marc Pr\'evost,
\emph{Remainder Pad\'e approximants for the Hurwitz zeta function},
Results in Mathematics \textbf{74} (2019), article 51.
\href{https://arxiv.org/pdf/1709.05389v1}{arXiv:1709.05389v1},
15 September 2017. Theorems 1--2, equations (3.1)--(3.4), and Section 4.1
give the positive representation and moments used here.

\bibitem{prevost1996}
Marc Pr\'evost,
\emph{A new proof of the irrationality of \(\zeta(2)\) and \(\zeta(3)\)
using Pad\'e approximants},
Journal of Computational and Applied Mathematics \textbf{67} (1996),
219--235.
\href{https://doi.org/10.1016/0377-0427(95)00019-4}{DOI: 10.1016/0377-0427(95)00019-4}.
Historical context; the present proof does not invoke an unproved
result from this paper.

\bibitem{prevost-rivoal}
Marc Pr\'evost and Tanguy Rivoal,
\emph{Diagonal convergence of the remainder Pad\'e approximants for the
Hurwitz zeta function},
Journal of Number Theory \textbf{222} (2021), 346--361.
\href{https://doi.org/10.1016/j.jnt.2020.10.019}{DOI: 10.1016/j.jnt.2020.10.019}.
Background on approximation; no quantitative bound from this paper is used.

\bibitem{zudilin}
Wadim Zudilin,
\emph{A determinantal approach to irrationality},
Constructive Approximation \textbf{45} (2017), 301--310.
\href{https://arxiv.org/pdf/1507.05697v2}{arXiv:1507.05697v2},
especially Proposition 2 and Section 2.

\bibitem{uvarov}
V.\,B. Uvarov,
\emph{The connection between systems of polynomials that are orthogonal
with respect to different distribution functions},
USSR Computational Mathematics and Mathematical Physics
\textbf{9}, no.\ 6 (1969), 25--36.
\href{https://doi.org/10.1016/0041-5553(69)90124-4}{DOI: 10.1016/0041-5553(69)90124-4}.
Historical attribution through Krattenthaler's discussion \cite{krattenthaler};
the original article was not used as a proof source.

\bibitem{krattenthaler}
Christian Krattenthaler,
\emph{A determinant identity for moments of orthogonal polynomials that
implies Uvarov's formula for the orthogonal polynomials of rationally
related densities},
\href{https://arxiv.org/pdf/2103.03969v1}{arXiv:2103.03969v1}, 2021.
Introduction and Theorem 1 provide the general setting for rational
modification of a measure. The theorem itself is not required below.

\bibitem{brown}
Francis Brown,
\emph{Mellin transforms, transfinite diameter and rational approximations
of integrals},
\href{https://arxiv.org/html/2604.20741v1}{arXiv:2604.20741v1}, 2026.
Sections 3.1--3.2 and 7.4--7.5 give related pairing and normalization
frameworks. This is contextual credit, not a claim that Brown states
the specific estimates proved here.

\bibitem{andreief}
C. Andr\'eief,
\emph{Note sur une relation entre les int\'egrales d\'efinies des produits
des fonctions}, M\'emoires de la Soci\'et\'e des Sciences physiques et
naturelles de Bordeaux, third series, \textbf{2} (1886), 1--14.
For the identity and publication history we consulted Peter J. Forrester,
\emph{Meet Andr\'eief, Bordeaux 1886, and Andreev, Kharkov 1882--83},
\href{https://arxiv.org/pdf/1806.10411v1}{arXiv:1806.10411v1}, 2018.
The identity is proved in the text.

\bibitem{dlmf}
NIST Digital Library of Mathematical Functions,
\href{https://dlmf.nist.gov/4.36.E3}{4.36.3} (partial fractions of \(\coth\)),
\href{https://dlmf.nist.gov/25.6.E2}{25.6.2} (Euler's even-zeta formula),
\href{https://dlmf.nist.gov/25.11.E29}{25.11.29} (Hermite's formula),
\href{https://dlmf.nist.gov/24.2}{24.2} (Bernoulli generating functions),
\href{https://dlmf.nist.gov/24.4}{24.4} (Bernoulli identities), and
\href{https://dlmf.nist.gov/24.10.i}{24.10(i)}
(von Staudt--Clausen). Consulted 14 November 2026.

\end{thebibliography}

\end{document}
